\chapter{Introduction}

\section{Motivation}
Text Classification (TC) is a foundational task in Natural Language Processing (NLP) \cite{Zangari2024}. While early work often framed TC as a multiclass problem, modern texts regularly contain multiple overlapping themes \cite{Hu2025, TidakeSane2018}. This motivates Multi-Label Classification (MLC), where instances can be assigned multiple labels simultaneously \cite{TidakeSane2018}. Beyond mere multiplicity, labels are frequently correlated \cite{TidakeSane2018, Huang2024}, so modeling inter-label dependencies is crucial.

A particularly compelling setting is hierarchical multi-label classification (HMLC), where labels are arranged in a predefined hierarchy (e.g., tree or DAG) from coarse to fine categories \cite{liu2023recentadvanceshierarchicalmultilabel}. This structure matches how humans reason about complex discourse. In the high-stakes context of online news and crisis communication, identifying nested narratives is both socially relevant and technically challenging. As a further motivation, this work aims to rigorously evaluate the classification capabilities of large language models (LLMs) on a difficult, nuanced HMLC task—one that requires fine-grained distinctions, hierarchical consistency, and evidence-grounded reasoning rather than shallow keyword matching.

\section{Problem Definition}
We ground our study in SemEval-2025 Task 10, Subtask 2 \cite{semeval2025task10}, a hierarchical narrative detection task for news. The objective is to assign each document multiple labels from a two-level taxonomy of narratives and subnarratives, i.e., mapping a text to a set of (narrative, subnarrative) pairs. The task features: (i) extreme class imbalance; (ii) hierarchical consistency constraints; and (iii) multilingual complexities. Formally, the MLC problem seeks to construct a function $f: X \to 2^L$ that maps an instance $x \in X$ to a subset $Y \subseteq L$, where $L$ is the finite label set \cite{TidakeSane2018}. This function may be learned through supervised training or constructed via other means (e.g., zero-shot prompting). HMLC adds the requirement that predicted labels respect the hierarchical structure \cite{liu2023recentadvanceshierarchicalmultilabel}.

\section{Related Work}
Research in multi-label classification has progressed from early problem-transformation methods (Binary Relevance, Classifier Chains, Label Powerset) \cite{zhang_binary_2018, weng_label_2020, shan_co-learning_2018} to Transformer-based models that learn contextual correlations \cite{devlin_bert_2019} and then to hierarchy-aware architectures incorporating structure through specialized losses or graph formulations \cite{sadat_hierarchical_2022, peng_hierarchical_2021, gong_hierarchical_2020, liu2023recentadvanceshierarchicalmultilabel}. Yet in settings that are simultaneously long-tailed, multilingual, and hierarchically constrained, supervised approaches remain brittle: they overfit head labels, struggle with rare-node calibration, propagate errors top–down, and provide limited transparency about individual decisions. Recent zero-shot and agentic LLM approaches \cite{GateNLP_github2025, Singh2025} reduce dependence on labeled data, but introduce new gaps in traceability, cost efficiency, and consistent enforcement of hierarchical constraints (true-path validity, evidence grounding). This trajectory motivates our contributions: systematically contrasting a strong supervised baseline with two zero-shot architectural paradigms while explicitly targeting the unresolved challenges of (i) long-tail robustness, (ii) native multilingual operation without translation-induced semantic loss, and (iii) auditable, evidence-grounded reasoning.

\section{Our Contributions}
We make the following contributions in the context of HMLC for narrative detection:
\begin{itemize}
  \item We introduce a zero-shot agentic framework that decomposes narrative detection into specialized binary decisions handled by LLM agents, achieving competitive performance without task-specific fine-tuning.
  \item We propose a configurable, traceable actor - critic pipeline that augments the agentic approach with explicit evidence extraction and validation, improving robustness and interpretability.
  \item We provide a comparative analysis of three paradigms: a fine-tuned Transformer baseline, the agentic framework, and the actor - critic pipeline, highlighting trade-offs in accuracy, traceability, and cost.
  \item We conduct an in-depth error analysis on SemEval-2025 data, elucidating the roles of class imbalance, semantic ambiguity, and hierarchical constraints in model failures.
\end{itemize}

\paragraph{Paper outline.} The remainder of this paper presents task and data analysis, details the system architectures, reports experiments and ablations, and concludes with discussion and future directions.
